\documentclass[troisieme]{fiche}
\metadonnees{13}{Un thé sans fin}{Bloc C — Décrire et comparer}

\begin{document}
\entete

\objectifs{%
\textbf{Langue} : dénombrables et indénombrables ; \emph{some} / \emph{any}, \emph{much} /
\emph{many}.\par
\textbf{Texte} : Lewis Carroll, \emph{Alice's Adventures in Wonderland} (1865),
ch. \textsc{vii}.\par
\textbf{Durée} : 1 h — exercices 20 min, texte et questions 25 min, version 15 min.}

%% ===================================================================
\exercice[12 min]{\emph{Some}, \emph{any}, \emph{much}, \emph{many}}

\rappel{\textbf{Rappel.} \emph{Some} à l'affirmative, \emph{any} à la question et à la
négative. \emph{Much} devant un indénombrable, \emph{many} devant un pluriel.}

\consigne{Complétez avec \emph{some}, \emph{any}, \emph{much}, \emph{many}, \emph{a few} ou
\emph{a little}.}

\begin{anglais}
\begin{enumerate}[label=\arabic*.,itemsep=2.5mm,leftmargin=*]
  \item There was nothing on the table but tea: there wasn't \trou{any} wine.
  \item \trou{Some} of the guests were fast asleep.
  \item How \trou{many} cups are there on the table?
  \item How \trou{much} tea did you drink?
  \item There isn't \trou{much} room at this corner.
  \item Alice asked \trou{a few} questions and got no answer.
  \item Would you like \trou{some} more bread and butter?
  \item He bought \trou{a few} things at the shop on the corner.
  \item There is only \trou{a little} butter left.
  \item Did they sell \trou{any} clothes in that street?
\end{enumerate}
\end{anglais}

\note[Les deux règles]{\emph{Some} / \emph{any} : la seule question qui compte est celle de la
phrase — affirmative (2), interrogative (10), négative (1). \emph{Much} / \emph{many} :
la seule question est celle du nom — \emph{tea} et \emph{room} ne se comptent pas (4, 5),
\emph{cups} se comptent (3).}

\note[Deux exceptions utiles]{Dans une offre ou une demande polie, on emploie \emph{some} et
non \emph{any} : \emph{would you like some tea?} (7). Et à l'affirmative, \emph{much} et
\emph{many} sont rares : on dit plutôt \emph{a lot of}.}

\note[\emph{A few} et \emph{a little}]{\emph{A few} + pluriel (6, 8), \emph{a little} +
indénombrable (9). Sans l'article, le sens se retourne : \emph{few friends} = presque pas
d'amis ; \emph{a few friends} = quelques amis.}

%% ===================================================================
\exercice[8 min]{Les indénombrables anglais}

\rappel{\textbf{Rappel.} \emph{News}, \emph{advice}, \emph{furniture}, \emph{hair},
\emph{information} sont indénombrables : jamais de pluriel, verbe au singulier.}

\consigne{Traduisez en anglais.}

\begin{anglais}
\begin{enumerate}[label=\arabic*.,itemsep=3mm,leftmargin=*]
  \item J'ai une bonne nouvelle.
    \lignes{1}
    \corr{I have some good news. \emph{ou} I have a piece of good news.}
  \item Tes cheveux sont trop longs.
    \lignes{1}
    \corr{Your hair is too long.}
  \item Il m'a donné un conseil.
    \lignes{1}
    \corr{He gave me a piece of advice.}
  \item Il n'y a pas beaucoup de meubles dans la pièce.
    \lignes{1}
    \corr{There isn't much furniture in the room.}
  \item Combien de pain reste-t-il ?
    \lignes{1}
    \corr{How much bread is left?}
  \item Je voudrais des renseignements.
    \lignes{1}
    \corr{I would like some information.}
\end{enumerate}
\end{anglais}

\note[Compter l'incomptable]{Pour isoler une unité, l'anglais passe par \emph{a piece of} :
\emph{a piece of news}, \emph{a piece of advice}, \emph{a piece of furniture}. On ne dit ni
\emph{an advice} ni \emph{advices}.}

\note[\emph{Hair}]{Singulier pour la chevelure : \emph{her hair is long}. Le pluriel
\emph{hairs} existe, mais il désigne les cheveux un à un — celui qu'on trouve dans la soupe.}

%% ===================================================================
\newpage
\section{Texte}

\noindent\small\textit{Alice arrive devant la maison du Lièvre de Mars. Une table est dressée
dehors.}\normalsize

\begin{texte}
There was a table set out under a tree in front of the house, and the March
\gl[a hare]{Hare}{un lièvre} and the \gl[a hatter]{Hatter}{un chapelier} were having tea at
it: a \gl[a dormouse]{Dormouse}{un loir} was sitting between them, fast asleep, and the other
two were using it as a \gl[a cushion]{cushion}{un coussin}, resting their
\gl[an elbow]{elbows}{un coude} on it, and talking over its head. ``Very uncomfortable for the
Dormouse,'' thought Alice; ``only, as it's asleep, I suppose it doesn't mind.''

The table was a large one, but the three were all \gl[to crowd]{crowded}{s'entasser} together
at one corner of it: ``No room! No room!'' they cried out when they saw Alice coming.
``There's \emph{plenty} of room!'' said Alice \gl[indignantly]{indignantly}{avec indignation},
and she sat down in a large arm-chair at one end of the table.

``Have some wine,'' the March Hare said in an encouraging tone.

Alice looked all round the table, but there was nothing on it but tea. ``I don't see any
wine,'' she remarked.

``There isn't any,'' said the March Hare.

``Then it wasn't very \gl[civil]{civil}{poli, courtois} of you to offer it,'' said Alice
angrily.

``It wasn't very civil of you to sit down without being invited,'' said the March Hare.

``I didn't know it was \emph{your} table,'' said Alice; ``it's \gl[to lay, laid]{laid}{mettre
(le couvert)} for a great many more than three.''

``Your hair wants cutting,'' said the Hatter. He had been looking at Alice for some time with
great curiosity, and this was his first speech.

``You should learn not to make personal remarks,'' Alice said with some
\gl[severity]{severity}{sévérité}; ``it's very \gl[rude]{rude}{impoli}.''

The Hatter opened his eyes very wide on hearing this; but all he \emph{said} was, ``Why is a
\gl[a raven]{raven}{un corbeau} like a \gl[a writing-desk]{writing-desk}{un secrétaire,
un bureau}?''

``Come, we shall have some fun now!'' thought Alice. ``I'm glad they've begun asking
\gl[a riddle]{riddles}{une devinette}. --- I believe I can guess that,'' she added
\gl[aloud]{aloud}{à voix haute}.

``Do you mean that you think you can find out the answer to it?'' said the March Hare.

``Exactly so,'' said Alice.

``Then you should say what you mean,'' the March Hare went on.

``I do,'' Alice \gl[hastily]{hastily}{précipitamment} replied; ``at least --- at least I mean
what I say --- that's the same thing, you know.''

``Not the same thing a bit!'' said the Hatter. ``You might just as well say that `I see what I
eat' is the same thing as `I eat what I see'!''
\end{texte}
\source{Lewis Carroll, \emph{Alice's Adventures in Wonderland}, Londres, Macmillan, 1865,
ch.~\textsc{vii}.}

\partiel[15 min]{Questions}
\consigne{Répondez aux deux premières questions en français, aux deux dernières en anglais.}

\begin{enumerate}[label=\arabic*.,itemsep=2.5mm,leftmargin=*]
  \item Pourquoi les trois convives crient-ils \og No room ! \fg{} alors que la table est
    grande ?
    \lignes{2}
    \corr{Parce qu'ils se sont tous entassés dans un seul coin de la table. Il y a de la place
    en abondance, comme Alice le fait remarquer : le manque de place est un prétexte pour
    écarter l'intruse.}
  \item Relevez trois impolitesses commises pendant la scène, et dites par qui.
    \lignes{3}
    \corr{Le Lièvre de Mars offre du vin qui n'existe pas ; il reproche à Alice de s'être
    assise sans y avoir été invitée ; le Chapelier lui dit que ses cheveux ont besoin d'être
    coupés. Alice, de son côté, s'assoit sans invitation et fait la leçon au Chapelier.}
  \item \en{Alice says she means what she says. Why do the Hatter and the March Hare disagree?}
    \lignes{3}
    \corr{Because the two sentences have the same words in a different order, and that changes
    everything: \emph{I see what I eat} is not \emph{I eat what I see}. Alice uses language as
    everybody does; they take it exactly as it is written.}
  \item \en{Find the two sentences with \emph{any} and the one with \emph{a great many}. Why
    \emph{any} and not \emph{some}?}
    \lignes{2}
    \corr{\emph{I don't see any wine} et \emph{There isn't any} : les deux phrases sont
    négatives. \emph{It's laid for a great many more than three} : \emph{many} parce que les
    couverts se comptent. On aurait \emph{some} dans une phrase affirmative — \emph{have some
    wine}, que le Lièvre dit justement au début.}
\end{enumerate}

\note[Une politesse renversée]{\emph{Your hair wants cutting} : \emph{to want} + \emph{-ing}
au sens de \og avoir besoin d'être \fg. Tournure encore courante en anglais britannique :
\emph{the car wants washing}.}

%% ===================================================================
\section{Version}
\consigne{Traduisez, de \emph{The Hatter opened his eyes very wide} à \emph{`I eat what I
see'}.}

\lignes{14}

\corr{\textbf{Proposition de traduction.} Le Chapelier ouvrit de grands yeux en entendant
cela ; mais tout ce qu'il \emph{dit} fut : \og Pourquoi un corbeau ressemble-t-il à un
secrétaire ? \fg{} \og Allons, nous allons nous amuser un peu ! pensa Alice. Je suis contente
qu'ils se soient mis à poser des devinettes. --- Je crois que je peux trouver celle-là \fg,
ajouta-t-elle à voix haute. \og Veux-tu dire que tu penses pouvoir en trouver la réponse ?
\fg{} demanda le Lièvre de Mars. \og Exactement \fg, dit Alice. \og Alors il faut dire ce que
tu penses \fg, reprit le Lièvre de Mars. \og Mais oui, répondit Alice précipitamment ; du
moins\ldots{} du moins je pense ce que je dis --- c'est la même chose, vous savez. \fg{} \og
Pas du tout la même chose ! dit le Chapelier. Autant dire que ``je vois ce que je mange''
est la même chose que ``je mange ce que je vois'' ! \fg}

\note[Choix de traduction 1]{\emph{all he said was} : \og tout ce qu'il dit fut \fg. Le
français a besoin de \og ce que \fg{} là où l'anglais se contente de \emph{all}.}

\note[Choix de traduction 2]{\emph{You might just as well say} : \og autant dire \fg. Locution
figée ; traduire \emph{might} par \og pourrais \fg{} donnerait une phrase illisible.}

\note[Choix de traduction 3]{Les guillemets dans les guillemets : la citation intérieure prend
des guillemets anglais, comme ci-dessus, car toute la réplique repose sur l'ordre exact des
deux phrases citées.}

%% ===================================================================
\partiel{Lexique à mémoriser}
\begin{multicols}{2}\small
\begin{itemize}[label=--,itemsep=0.8mm,leftmargin=*]
  \item \textbf{a hare} — un lièvre
  \item \textbf{a cushion} — un coussin
  \item \textbf{an elbow} — un coude
  \item \textbf{asleep} — endormi ; \textit{fast asleep} profondément endormi
  \item \textbf{room} — la place \textit{(indénombrable)}
  \item \textbf{plenty of} — beaucoup de, en abondance
  \item \textbf{rude} — impoli
  \item \textbf{a riddle} — une devinette
  \item \textbf{to guess} — deviner
  \item \textbf{to mean} — vouloir dire ; meant, meant
  \item \textbf{aloud} — à voix haute
  \item \textbf{hair} — les cheveux \textit{(indénombrable)}
\end{itemize}
\end{multicols}

\end{document}
